\section{Hyperelliptic specialization}
\label{sec:BA-realization}

The general construction becomes especially explicit on an odd-degree
hyperelliptic curve. Sheet reflection identifies ordinary and dual
functions, the kernel becomes a two-function Wronskian, and the residue
conditions become polynomial interpolation moments. We give the sigma
normalization and the conversion to \cref{thm:general-BA-realization},
then retain the explicit weights needed for the Mumford correspondence.
The mode counts and Hamiltonian remain
\cref{eq:BA-mode-counts,eq:BA-H0-flat}.

\subsection{Curve and function conventions}

Fix a smooth odd-degree curve, its Abel map and ordered differentials:
\begin{equation}\label{eq:BA-curve}
 y^2=f(x)=4x^{2g+1}+\sum_{r=0}^{2g}\lambda_r x^r,\qquad
 A(P)=\int_\infty^P du,\qquad
 du_j=\frac{x^{j-1}\,dx}{y},\quad1\leq j\leq g.
\end{equation}
The sheet involution is $P=(a,y)\mapsto\iota(P)=(a,-y)$;
compatible lifts give $A(\iota(P))=-A(P)$.
This is Generapy's automatic odd-degree Baker basis \cite{Generapy}.
The vector $z$ in this section uses these Baker Abel coordinates,
related by a fixed linear change to the normalized theta coordinates
of \cref{eq:general-BA-theta}. The scalar time follows the same first BA direction, with first BA
parameter $-it$.
In converting
from a monic equation $y_{\mathrm{mon}}^2=f_{\mathrm{mon}}(x)$ with
differentials $x^{j-1}dx/(2y_{\mathrm{mon}})$, set
$y=2y_{\mathrm{mon}}$ and $f=4f_{\mathrm{mon}}$; the Abel coordinates
then agree.

Use the canonical hyperelliptic Baker sigma function, with
$\wp_{ij}=-\partial_{z_i}\partial_{z_j}\log\sigma$ and
\begin{equation}\label{eq:BA-sharp}
 \sigma_\sharp=\sigma_{2,4,6,\ldots},
\end{equation}
where indices stop at $g$ and an empty list means $\sigma$ itself
\cite{BEL1997,Onishi2005}. To fix the multiplicative sigma constant,
its lowest ordinary-degree term is
$\det(z_{i+j-1})_{1\leq i,j\leq\lfloor(g+1)/2\rfloor}$.
These are the normalization and derivative order used by Generapy's
\texttt{kleinian\_sigma} and \texttt{kleinian\_baker\_akhiezer}.

The compatible second-kind vector is
\begin{equation}\label{eq:BA-second-kind}
 \boldsymbol R(P)=\operatorname{FP}\int_\infty^P dr,\qquad
 dr_j=\frac1{4y}\sum_{k=j}^{2g-j}
          (k+1-j)\lambda_{k+1+j}x^k\,dx,
 \quad1\leq j\leq g,\qquad \lambda_{2g+1}=4.
\end{equation}
The finite part subtracts the polar primitive at infinity and fixes
its constant term to zero in the local parameter $x^{-1/2}$.
The first- and second-kind lifts use the same paths. With this
choice, second-kind half-periods satisfy
$2\eta=-\oint_a dr$, and the exponential in the BA function has
the minus sign
\begin{equation}\label{eq:BA-generapy}
 \Psi(P,z)=\frac{\sigma(A(P)-z)}
 {\sigma_\sharp(A(P))\sigma(z)}
 \exp\bigl(-\boldsymbol R(P)^{\trans}z\bigr).
\end{equation}
This is the Generapy convention; see also the second-kind and
Schr\"odinger conventions in \cite{Christiansen2000} and the
normalized genus-two formula in \cite[equation (3.3)]{Braden2005}.

The physical time follows the last Abel direction:
\begin{equation}\label{eq:BA-time-polynomial}
 z(t)=z_0-it e_g,\qquad
 \mathcal P(z,x)=x^g-\sum_{j=1}^{g}\wp_{gj}(z)x^{j-1},
 \qquad \mathcal P_{z_g}=\partial_{z_g}\mathcal P=i\dot{\mathcal P}.
\end{equation}
Here $z_0$ is unrestricted and $e_g$ is the last coordinate vector.
The polynomial $f(x)$ fixes the curve and has degree $2g+1$;
$\mathcal P(z,x)$ has degree $g$ and describes the moving divisor.

The standard one-point identities, in these conventions, are
\begin{equation}\label{eq:BA-identities}
\begin{aligned}
 \Psi(P,z)\Psi(\iota(P),z)&=(-1)^g\mathcal P(z,a),\\
 \partial_{z_g}\log\Psi(P,z)
 &=\frac{\mathcal P_{z_g}(z,a)-y}{2\mathcal P(z,a)},\\
 \partial_{z_g}^2\Psi(P,z)
 &=\left(a+2\wp_{gg}(z)+\frac{\lambda_{2g}}4\right)\Psi(P,z).
\end{aligned}
\end{equation}
They hold for a free $z$ and a single curve-point shift $A(P)$
\cite{BEL1997,Eilbeck2000}. In particular, raw sheet partners have
the sign $(-1)^g$ in their product. We absorb that constant sign
into the mode normalization so that their power is a weighted
evaluation of the monic $\mathcal P$.

\subsection{Sigma modes and the general normalization}

Use the mode indices and canonical partners of
\cref{eq:BA-mode-dictionary}. In this section the constants
$\kappa_k$ multiply the raw sigma products in
\cref{eq:BA-powers-rates}; their conversion to the normalized
residue weights of \cref{eq:general-residue-weights} is given in
\cref{eq:BA-general-dictionary}. The involution $\iota$ denotes sheet reflection; bars label independent Hamiltonian partners.

Choose $K$ points $Q_k=(a_k,y_k)$ on \cref{eq:BA-curve} with
finite distinct abscissae, and nonzero constants $\kappa_k,c_k$,
$1\leq k\leq K$. Define all elementary modes by
\begin{equation}\label{eq:BA-elementary}
 \xi_k=c_k\Psi(Q_k,z),\qquad
 \bar{\xi}_k=\frac{(-1)^g\kappa_k}{c_k}\Psi(\iota(Q_k),z),
 \qquad1\leq k\leq K.
\end{equation}
Their powers and logarithmic rates are, for $1\leq k\leq K$,
\begin{equation}\label{eq:BA-powers-rates}
\begin{aligned}
 \xi_k\bar{\xi}_k&=\kappa_k\mathcal P(z,a_k),\\
 \frac{i\dot\xi_k}{\xi_k}
 &=\frac{\mathcal P_{z_g}(z,a_k)-y_k}{2\mathcal P(z,a_k)},&
 \frac{i\dot{\bar{\xi}}_k}{\bar{\xi}_k}
 &=\frac{\mathcal P_{z_g}(z,a_k)+y_k}{2\mathcal P(z,a_k)}.
\end{aligned}
\end{equation}
Define the Wronskian in the Abel direction by
$W(f_1,f_2)=f_1\partial_{z_g}f_2-(\partial_{z_g}f_1)f_2$
for arbitrary functions $f_1,f_2$ of $z$. Along the physical flow,
$W(f_1,f_2)=i(f_1\dot f_2-\dot f_1f_2)$.
Each elementary pair has the constant Abel-direction Wronskian
$W(\xi_k,\bar\xi_k)=\kappa_k y_k$, $1\leq k\leq K$.
All elementary modes share the potential
$2\wp_{gg}+\lambda_{2g}/4$ in \cref{eq:BA-identities}.

For every cross-group pair set
\begin{equation}\label{eq:BA-pair-composites}
\begin{aligned}
 \xi_{ij}
 &=\frac{\sqrt2}{a_j-a_i}
        W(\bar{\xi}_i,\bar{\xi}_j),\\
 \bar{\xi}_{ij}
 &=-\frac{\sqrt2}{a_j-a_i}
        W(\xi_i,\xi_j),
 \qquad 1\leq i\leq L,\quad L+1\leq j\leq K.
\end{aligned}
\end{equation}
The kernel specialization in \cref{eq:BA-hyperelliptic-kernel} identifies these Wronskians with
sigma ratios having shifts $A(Q_i)+A(Q_j)$ and its negative,
with fixed normalizing factors
\cite[Theorem~6.2, $n=2$]{Nakayashiki2026}. Their differential equations have
a particularly direct explanation: the common potential cancels,
leaving
\begin{equation}\label{eq:BA-pair-rates}
 i\dot\xi_{ij}=\sqrt2\bar{\xi}_i\bar{\xi}_j,
 \qquad
 i\dot{\bar{\xi}}_{ij}=-\sqrt2\xi_i\xi_j,
 \qquad1\leq i\leq L,\quad L+1\leq j\leq K.
\end{equation}
These are exactly the pair-mode equations of $H_0$. The elementary equations are supplied by the general realization criterion.
Primitives of products of ordinary and dual BA functions also occur in
the Cauchy--BA kernel construction
\cite[Section~2.4, equation~(2.13)]{GrinevichOrlov1989} and in
squared-eigenfunction-potential theory
\cite[Corollary~3.1]{AratynNissimovPacheva1997}.
The residue criterion takes the explicit polynomial form developed in the next subsection.

Here is the explicit dictionary. Put $\epsilon=(-1)^g$ and
$P_0(x)=\mathcal P(z_0,x)$. The ordinary and dual functions of
\cref{eq:general-BA-asymptotics} are
\begin{equation}\label{eq:BA-normalized-sigma}
 \psi(P,t)=\frac{\epsilon\Psi(\iota(P),z(t))}{\epsilon\Psi(\iota(P),z_0)},
 \qquad
 \psi^\vee(P,t)=\frac{\Psi(P,z(t))}{\Psi(P,z_0)},\qquad
 d\Omega_0=P_0(x)\frac{dx}{y}.
\end{equation}
The raw product identity gives $\psi\psi^\vee=\mathcal P/P_0$. For a regular
initial chart $P_0(a_k)\ne0$, take the general construction's points,
weights and gauges to be
\begin{equation}\label{eq:BA-general-dictionary}
\begin{aligned}
 \widehat Q_i&=Q_i,&&1\leq i\leq L,&
 \widehat Q_j&=\iota(Q_j),&&L+1\leq j\leq K,\\
 \widehat\kappa_k&=\kappa_kP_0(a_k),&&1\leq k\leq K,&
 \widehat c_k&=c_k\Psi(Q_k,z_0),&&1\leq k\leq K.
\end{aligned}
\end{equation}
The hats identify the normalized constants used in
\cref{eq:general-elementary-modes,eq:general-pair-modes}.
The kernel is
\begin{equation}\label{eq:BA-hyperelliptic-kernel}
 \Omega(P,Q)=\frac{\psi^\vee(Q)\partial_{z_g}\psi(P)-\partial_{z_g}\psi^\vee(Q)\psi(P)}{x(P)-x(Q)}.
\end{equation}
Its Abel-direction derivative is $\psi(P)\psi^\vee(Q)$ by the common Schr\"odinger equation.
Substituting \cref{eq:BA-general-dictionary} into
\cref{eq:general-pair-modes} gives exactly the Wronskians in
\cref{eq:BA-pair-composites}, including their signs and $\sqrt2$
factors. Thus the elementary and pair functions are the same modes,
expressed in normalized BA and raw sigma conventions.

\subsection{Explicit interpolation weights and phases}
\label{sec:BA-interpolation}

In the local parameter $s=x^{-1/2}$ at infinity,
$du_\alpha=-s^{2g-2\alpha}(1+O(s^2))ds$ for
$1\leq\alpha\leq g$. At the finite nodes, use $x-a_k$.
The general residue conditions therefore become
\begin{equation}\label{eq:BA-moments-final}
 \sum_{i=1}^{L}\kappa_i a_i^r
 =\sum_{j=L+1}^{K}\kappa_j a_j^r
 =\begin{cases}
 0,&0\leq r\leq g-2,\\
 \tfrac12,&r=g-1.
 \end{cases}
\end{equation}
Reflection changes both the principal coefficient of the interpolation
function and the differential coefficient, leaving identical moments
for the two groups.

To construct all weights with these moments, collect their simple
fractions over the node polynomials
\begin{equation}\label{eq:BA-node-polynomials}
 R_L(x)=\prod_{i=1}^{L}(x-a_i),\qquad
 R_M(x)=\prod_{j=L+1}^{K}(x-a_j).
\end{equation}
Define their numerator polynomials $q_L,q_M$ by
\begin{equation}\label{eq:BA-weight-generating}
 \frac{q_L(x)}{R_L(x)}
 =2\sum_{i=1}^{L}\frac{\kappa_i}{x-a_i},\qquad
 \frac{q_M(x)}{R_M(x)}
 =2\sum_{j=L+1}^{K}\frac{\kappa_j}{x-a_j}.
\end{equation}
Before imposing the moments, these numerators have degrees at most
$L-1,M-1$. The moments say precisely that each rational function has expansion
$x^{-g}+O(x^{-g-1})$ at infinity. Hence $L,M\geq g$, and the
numerators must be monic of degrees $L-g,M-g$. Taking residues gives
\begin{equation}\label{eq:BA-residue-weights}
 \kappa_i=\frac{q_L(a_i)}{2R_L'(a_i)},
 \quad1\leq i\leq L,\qquad
 \kappa_j=\frac{q_M(a_j)}{2R_M'(a_j)},
 \quad L+1\leq j\leq K.
\end{equation}
Primes denote differentiation in $x$. Conversely, every pair of
monic numerators of these degrees gives the required moments;
their values at the respective nodes must be nonzero to keep the
mode weights nonzero. Thus $q_L,q_M$
parametrize the available interpolation weights: they equal $1$
when the corresponding group has exactly $g$ points, while larger
groups give free polynomial coefficients.

\begin{corollary}\label{thm:BA-realization}
Let the elementary modes and pair composites in
\cref{eq:BA-elementary,eq:BA-pair-composites} range over a nonempty
open subset of the Jacobian of the fixed curve
\cref{eq:BA-curve}. They satisfy all Hamilton equations of
\cref{eq:BA-H0-flat} if and only if their weights obey the interpolation
moments \cref{eq:BA-moments-final} and their phase coefficients satisfy
\begin{equation}\label{eq:BA-phase-final}
\begin{aligned}
 b_i
 &=-\sum_{j=L+1}^{K}\kappa_j
       \frac{y_j-y_i}{a_j-a_i},&&1\leq i\leq L,\\
 b_j
 &=-\sum_{i=1}^{L}\kappa_i
       \frac{y_i-y_j}{a_i-a_j},&&L+1\leq j\leq K.
\end{aligned}
\end{equation}
\end{corollary}
\begin{proof}
In the unnormalized hyperelliptic weights the two interpolation
functions are
\begin{equation}\label{eq:BA-hyperelliptic-h}
 h_L(P)=\sum_{i=1}^{L}\kappa_i\frac{y(P)+y_i}{x(P)-a_i},\qquad
 h_M(P)=\sum_{j=L+1}^{K}\kappa_j\frac{y(P)-y_j}{x(P)-a_j}.
\end{equation}
Their finite poles are at $Q_i$ and $\iota(Q_j)$, respectively.
The moments make their puncture expansion $s^{-1}+O(s)$.
With \cref{eq:BA-normalized-sigma}, half the residue of
$h\,d\Omega_0$ is $\kappa_kP_0(a_k)$, precisely the normalized
weight in \cref{eq:BA-general-dictionary}.
Their opposite-group values give \cref{eq:BA-phase-final}.
The equivalence now follows from \cref{thm:general-BA-realization}.
At a branch point use its regular ordinate coordinate for the residue;
the formula extends without division by its zero ordinate.
\end{proof}

The phase conditions \cref{eq:BA-phase-final} explain the role of the
independent constant diagonal blocks. For $N=1$, the printed
Jiang--Rauch-Wojciechowski specialization gives
$b_k=b_1^N$ for every $1\leq k\leq K$, whereas these conditions
generically require different rates at different marked points.
The outer diagonal parameters in \cref{eq:BA-phase-dictionary}
supply exactly that freedom. The constants $c_k$ specify initial mode
normalizations and leave these logarithmic rates unchanged.

For example, take $g=L=M=2$ and let $Y(x)$ be the polynomial of degree
at most three with $Y(a_k)=y_k$, $1\leq k\leq4$. The moment weights
\cref{eq:BA-residue-weights} give
\begin{equation}\label{eq:BA-phase-differences}
 b_1-b_2=-\tfrac12[x^3]Y(x)(a_1-a_2),\qquad
 b_3-b_4=-\tfrac12[x^3]Y(x)(a_3-a_4).
\end{equation}
Thus a nonzero cubic interpolation coefficient requires distinct
phase rates within both groups. This is the generic minimal genus-two
realization. \Cref{thm:canonical-independence} applies both to these
distinct rates and to their coincident specializations.

\Cref{prop:BA-prescribed-spectral-curve} identifies this hyperelliptic
family as a spectral level of genus $g$ with total defect $LM-g$.

\subsection{Half-period shifts and al functions}
\label{sec:BA-al-functions}

Branch-point evaluations give a classical specialization of these
canonical modes. Let $B_k=(a_k,0)$ be a finite branch point and
$\delta_k=A(B_k)$. Since
$\operatorname{div}(x-a_k)=2B_k-2\infty$, Abel's theorem gives
$2\delta_k\in\Lambda$, where
$\Lambda=2\omega\mathbb Z^g+2\omega'\mathbb Z^g$.
The two sheet-partner functions are then proportional, and their
Wronskian vanishes. After a constant normalization they are
Weierstrass's hyperelliptic al functions, satisfying
\begin{equation}\label{eq:BA-al-power}
 \operatorname{al}_k(z)=d_k\Psi(B_k,z),\qquad
 \operatorname{al}_k(z)^2=\mathcal P(z,a_k),\qquad d_k^2=(-1)^g.
\end{equation}
Each al function is meromorphic on a double cover of the Jacobian;
its square is Abelian. These functions generalize the Jacobi elliptic
functions \cite[Appendix]{MatsutaniPreviato2015}.

For an arbitrary-genus example containing all $2g+1$ finite branch
points, partition them into groups of sizes $L=g$, $M=g+1$.
Choose $q_L=1$ and $q_M(x)=x-\alpha$, with $\alpha$ distinct from
the second group's abscissae, and use the weights in
\cref{eq:BA-residue-weights}. They satisfy
\cref{eq:BA-moments-final}. All ordinates are zero, so
\cref{eq:BA-phase-final} gives $b_k=0$, $1\leq k\leq2g+1$.
Thus \cref{thm:BA-realization} realizes the entire elementary al
family in the purely cubic Hamiltonian \cref{eq:BA-H0-flat}.
The cross-group pair modes are their Wronskians in
\cref{eq:BA-pair-composites}, with half-period sums as shifts.
The canonical partners remain independent in the ambient system;
they become proportional on this al solution family.

Since all phases vanish, this family also solves the original
Jiang--Rauch-Wojciechowski system of
Appendix~\ref{sec:JRW-dictionary}; its complete integral count is
provided by the repeated-phase part of
\cref{thm:canonical-independence}, with moment integrals from blocks
of sizes $g$ and $g+1$.
In the product-surface curve $\Gamma$ of
\cref{prop:BA-prescribed-spectral-curve}, these groups meet the boundary
at two points of multiplicities $g$ and $g+1$.
Their defects are at least $g(g-1)/2$ and $g(g+1)/2$,
whose sum $g^2=LM-g$ accounts for the entire defect.
Consequently the affine curve is smooth.

For genus at least two, the elementary half-period mechanism is
specific to hyperelliptic curves: $2A(P)\in\Lambda$ for $P\ne P_\infty$ makes
$2P-2P_\infty$ principal and supplies a degree-two map.
Other covering structures give higher-torsion analogues. On a cyclic
trigonal curve $y^3=\prod_{r=1}^{4}(x-b_r)$, for example,
$\operatorname{div}(x-b_r)=3B_r-3\infty$ gives order-three shifts;
trigonal al functions are developed in
\cite[Section~8]{MatsutaniPreviato2015}.

\subsection{Genus one and many elementary modes}

In Weierstrass coordinates take
$y^2=4x^3-g_2x-g_3$. Then $A(P)=\mu$,
$a=\wp(\mu)$, $y=\wp'(\mu)$ and
$R(P)=-\zeta(\mu)$, so the elementary function is
\begin{equation}\label{eq:BA-elliptic}
 \Psi(P,z)=\frac{\sigma(\mu-z)}{\sigma(\mu)\sigma(z)}
                  e^{\zeta(\mu)z},\qquad z=z_0-it.
\end{equation}
Its normalized mode powers are
$\xi_k\bar{\xi}_k=\kappa_k(a_k-\wp(z))$, $1\leq k\leq K$.
The only moment conditions are
\[
 \sum_{i=1}^{L}\kappa_i
 =\sum_{j=L+1}^{K}\kappa_j=\tfrac12.
\]
Thus arbitrarily large groups of elementary elliptic modes fit the
same cubic-interaction construction, with their $LM$ pair composites.

For $L=M=1$, so $K=2$, write
$Y=(y_2-y_1)/(2(a_2-a_1))$.
The two weights are $1/2$, and phase matching gives
$b_1=b_2=-Y$. The resulting three-pair Hamiltonian is
\begin{equation}\label{eq:BA-three-wave}
 H_0=-Y(\xi_1\bar{\xi}_1+\xi_2\bar{\xi}_2)
 +\sqrt2(\xi_1\xi_2\xi_{12}
             +\bar{\xi}_1\bar{\xi}_2\bar{\xi}_{12}).
\end{equation}
This recovers the familiar detuned three-wave form
\cite{MenyukChenLee1983Integrable}. Increasing $L,M$ retains elliptic
elementary solutions while adding pair-composite interactions; it
also increases the genus of the regular ambient spectral curve.
